% !TeX root = ../tesis.tex


\chapter{Eritrocito sano}

\label{section:apendix3}

En la Sección \ref{section:Res_Eritrocito_Sano} se discutieron los resultados obtenidos de $Q_{\text{ext}}$ y del patrón de esparcimiento para $\lambda=431$ nm, correspondiente a la banda de Soret. En esta sección se presentan los resultados para las longitudes de onda representativas de las bandas $\beta$ y $\alpha$. En la Fig.~\ref{subfig:FF_541_health} se muestra la sección transversal diferencial de esparcimiento $\dd{C_{\text{sca}}}/\dd{\Omega}$, así como su representación polar [Fig. \ref{subfig:FF_polar_541_health}] para los cortes $\varphi=0^{\circ}$ y $\varphi=90^{\circ}$ para $\lambda=541$ nm. En la Fig. \ref{fig:FF_cartesian_541_health} se muestra la representación cartesiana de $\dd{C_{\text{sca}}}/\dd{\Omega}$ para los cortes $\varphi=0^{\circ}$ y $\varphi=90^{\circ}$ para $\lambda=541$ nm. Al igual que para $\lambda=431$ nm, el esparcimiento se concentra predominantemente en la región frontal, con un factor de anisotropía de $g=0.98$. Por otra parte, la diferencia entre los cortes $\varphi=0^{\circ}$ y $\varphi=90^{\circ}$ continúa localizándose principalmente en la región lateral, donde la correlación disminuye a $r=0.49$. 
	%
\begin{figure}[h!]
	\centering
	\sidesubfloat[First image]{\hspace{-0.5cm}{	\includegraphics[width=7.3cm]{../../Figuras/4-Resultados/FFdistribution541_EriSano.pdf}\label{subfig:FF_541_health}}}\hspace{0.1cm}
	\sidesubfloat[Second image]{\hspace{-0.5cm}{
			\includegraphics[width=8cm]{../../Figuras/4-Resultados/polarplot541FF_EriSano.pdf}\label{subfig:FF_polar_541_health}}}\vspace{-0.1cm}
	\caption{\textbf{a)} Sección transversal diferencial de esparcimiento $\dd{C_{\text{sca}}}/\dd{\Omega}$ normalizada respecto a su valor máximo para una longitud de onda $\lambda=541$ nm en función del ángulo de esparcimiento $\theta$ y del ángulo $\varphi$ para un eritrocito sano. Se identifican dos planos $\varphi=0^{\circ}$ en azul claro y $\varphi=90^{\circ}$ en azul oscuro. \textbf{b)} Sección transversal diferencial de esparcimiento $\dd{C_{\text{sca}}}/\dd{\Omega}$ en decibeles en función de $\theta$ para dos cortes: $\varphi=0^{\circ}$ (puntos azules claro) y $\varphi=90^{\circ}$ (puntos azules oscuro). La línea continua que une los puntos es una guía visual.}
	\label{fig:FF_541_health_1}
\end{figure}
%

	%
\begin{figure}[h!]
	\centering
			\includegraphics[width=14cm]{../../Figuras/4-Resultados/linearplot541FF_EriSano.pdf}\vspace{-0.1cm}
	\caption{Sección transversal diferencial de esparcimiento $\dd{C_{\text{sca}}}/\dd{\Omega}$ en función de $\theta$ para dos cortes: $\varphi=0^{\circ}$ (puntos azules claro) y $\varphi=90^{\circ}$ (puntos azules oscuro). La línea continua que une los puntos es una guía visual.}
	\label{fig:FF_cartesian_541_health}
\end{figure}
%
En la Fig.~\ref{subfig:FF_573_health} se muestra la sección transversal diferencial de esparcimiento $\dd{C_{\text{sca}}}/\dd{\Omega}$, así como su representación polar [Fig. \ref{subfig:FF_polar_573_health}] para los cortes $\varphi=0^{\circ}$ y $\varphi=90^{\circ}$ para $\lambda=573$~nm. En la Fig. \ref{fig:FF_cartesian_573_health} se muestra la representación cartesiana de $\dd{C_{\text{sca}}}/\dd{\Omega}$ para los cortes $\varphi=0^{\circ}$ y $\varphi=90^{\circ}$ para $\lambda=573$ nm. Al igual que para las longitudes de onda de 431 y 541 nm, el esparcimiento se concentra predominantemente en la región frontal. En este caso, el factor de anisotropía alcanza su valor máximo ($g=0.99$), lo que indica un mayor esparcimiento hacia adelante. Asimismo, las principales diferencias entre los cortes $\varphi=0^{\circ}$ y $\varphi=90^{\circ}$ continúan localizándose en la región lateral, donde la correlación aumenta a $r=0.71$, valor superior al obtenido para $\lambda=541$ nm, aunque todavía inferior al observado en las regiones frontal y de retroesparcimiento. En conjunto, los resultados para las tres longitudes de onda muestran que el esparcimiento permanece predominantemente dirigido hacia adelante y que dicha direccionalidad aumenta para las longitudes de onda consideradas, como se refleja en el incremento de $g$ de 0.97 a 0.99 y en las contribuciones integradas de la región frontal resumidas en la Tabla~\ref{table:431_val_health}. En contraste, la similitud entre los patrones laterales de ambos planos presenta un comportamiento no monótono con la longitud de onda, con un mínimo de correlación en $\lambda=541$ nm.
	%
\begin{figure}[tb]
	\centering
	\sidesubfloat[First image]{\hspace{-0.5cm}{	\includegraphics[width=7.3cm]{../../Figuras/4-Resultados/FFdistribution573_EriSano.pdf}\label{subfig:FF_573_health}}}\hspace{0.1cm}
	\sidesubfloat[Second image]{\hspace{-0.5cm}{
			\includegraphics[width=8cm]{../../Figuras/4-Resultados/polarplot573FF_EriSano.pdf}\label{subfig:FF_polar_573_health}}}\vspace{-0.1cm}
	\
	\vspace{0.1cm}
	\caption{\textbf{a)} Sección transversal diferencial de esparcimiento $\dd{C_{\text{sca}}}/\dd{\Omega}$ normalizada respecto a su valor máximo para una longitud de onda $\lambda=573$ nm en función del ángulo de esparcimiento $\theta$ y del ángulo $\varphi$ para un eritrocito sano. Se identifican dos planos $\varphi=0^{\circ}$ en azul claro y $\varphi=90^{\circ}$ en azul oscuro. \textbf{b)} Sección transversal diferencial de esparcimiento $\dd{C_{\text{sca}}}/\dd{\Omega}$ en decibeles en función de $\theta$ para dos cortes: $\varphi=0^{\circ}$ (puntos azules claro) y $\varphi=90^{\circ}$ (puntos azules oscuro). La línea continua que une los puntos es una guía visual.}
	\label{fig:FF_573_health_1}
\end{figure}
%

	%
\begin{figure}[h!]
	\centering
	
			\includegraphics[width=14cm]{../../Figuras/4-Resultados/linearplot573FF_EriSano.pdf}\vspace{-0.1cm}
	\caption{Sección transversal diferencial de esparcimiento $\dd{C_{\text{sca}}}/\dd{\Omega}$ en función de $\theta$ para dos cortes: $\varphi=0^{\circ}$ (puntos azules claro) y $\varphi=90^{\circ}$ (puntos azules oscuro). La línea continua que une los puntos es una guía visual.}
	\label{fig:FF_cartesian_573_health}
\end{figure}
%